\begin{table}[t]\centering\caption{Dyna-Q+ ablation, post-change reward, $n=10$, 20 seeds, mean (standard deviation). ``w/o both'' is Dyna-Q.}\label{tab:abl}
\resizebox{\columnwidth}{!}{\begin{tabular}{lccccc}
\toprule
Variant & blocking & shortcut & static & stoch. & stoch.\,blk \\
\midrule
Dyna-Q+ & 51.9 (4.1) & 229.7 (1.9) & 125.0 (0.8) & 36.5 (1.6) & 18.4 (1.2) \\
Dyna-Q+ w/o bonus & 36.8 (6.4) & 168.2 (3.4) & 124.3 (6.3) & 50.8 (3.5) & 13.1 (2.5) \\
Dyna-Q+ w/o untried & 17.3 (5.2) & 138.8 (18.2) & 52.1 (13.0) & 20.1 (3.7) & 3.1 (1.4) \\
Dyna-Q+ w/o both & 28.2 (8.4) & 165.8 (0.4) & 128.4 (3.9) & 46.1 (3.8) & 8.2 (2.7) \\
\bottomrule
\end{tabular}
}
\end{table}

\begin{table}[t]\centering\caption{Sensitivity of Dyna-Q+ ($n=10$, 10 seeds, mean) to the bonus weight $\kappa$; the last column is Dyna-Q on the same seeds.}\label{tab:kappa}
\resizebox{\columnwidth}{!}{\begin{tabular}{lrrrrr}
\toprule
Task (metric) & $10^{-4}$ & $10^{-3}$ & $10^{-2}$ & $3{\cdot}10^{-2}$ & Dyna-Q \\
\midrule
blocking (post) & 39.4 & 51.7 & 2.4 & 2.1 & 21.0 \\
shortcut (post) & 166.1 & 229.7 & 3.7 & 3.0 & 165.8 \\
static (cum) & 218.8 & 206.0 & 4.6 & 4.7 & 190.2 \\
stochastic (cum) & 86.7 & 66.9 & 9.2 & 9.8 & 78.4 \\
\bottomrule
\end{tabular}
}
\end{table}

\textbf{Which part of Dyna-Q+ matters (Table~\ref{tab:abl}).} Removing the bonus lowers post-change reward on the shortcut maze (\rhval{abl_dynaq_plus/dyna-q+-w-o-bonus/shortcut/post_reward/mean:1} versus \rhval{abl_dynaq_plus/dyna-q+/shortcut/post_reward/mean:1}, paired $p=\rhval{cmp/abl_dynaq_plus/dyna-q+-w-o-bonus/shortcut/post_reward/paired_p:sci1}$, \texttt{compare\_abl\_dynaq\_plus\_post\_reward.csv}), so the stale-transition bonus accounts for the shortcut gain, and the bonus-free variant equals Dyna-Q there within noise (\rhval{abl_dynaq_plus/dyna-q+-w-o-bonus/shortcut/post_reward/mean:1} versus \rhval{abl_dynaq_plus/dyna-q+-w-o-both/shortcut/post_reward/mean:1}). On blocking the drop (\rhval{abl_dynaq_plus/dyna-q+-w-o-bonus/blocking/post_reward/mean:1} versus \rhval{abl_dynaq_plus/dyna-q+/blocking/post_reward/mean:1}) is only a trend (paired $p=\rhval{cmp/abl_dynaq_plus/dyna-q+-w-o-bonus/blocking/post_reward/paired_p:3}$). Removing only the untried-action initialisation gives the lowest mean on every task: with the bonus but without that initialisation the agent is worse than plain Dyna-Q on the static maze (\rhval{abl_dynaq_plus/dyna-q+-w-o-untried/static/post_reward/mean:1} versus \rhval{abl_dynaq_plus/dyna-q+-w-o-both/static/post_reward/mean:1}, paired $p=\rhval{cmp/abl_ref_dynaq/dyna-q+-w-o-untried/static/post_reward/paired_p:sci1}$) and the stochastic maze (\rhval{abl_dynaq_plus/dyna-q+-w-o-untried/stochastic/post_reward/mean:1} versus \rhval{abl_dynaq_plus/dyna-q+-w-o-both/stochastic/post_reward/mean:1}, paired $p=\rhval{cmp/abl_ref_dynaq/dyna-q+-w-o-untried/stochastic/post_reward/paired_p:sci1}$; \texttt{compare\_abl\_ref\_dynaq\_post\_reward.csv}); on blocking, shortcut and stochastic-blocking the gaps to Dyna-Q are trends only ($p=\rhval{cmp/abl_ref_dynaq/dyna-q+-w-o-untried/blocking/post_reward/paired_p:2}$, $\rhval{cmp/abl_ref_dynaq/dyna-q+-w-o-untried/shortcut/post_reward/paired_p:2}$, $\rhval{cmp/abl_ref_dynaq/dyna-q+-w-o-untried/stoch_blocking/post_reward/paired_p:2}$). The two components are therefore not separable improvements. We have not run the interaction in a finer design. On the stochastic maze the bonus costs post-change reward (\rhval{abl_dynaq_plus/dyna-q+/stochastic/post_reward/mean:1} versus \rhval{abl_dynaq_plus/dyna-q+-w-o-bonus/stochastic/post_reward/mean:1} without it, paired $p=\rhval{cmp/abl_dynaq_plus/dyna-q+-w-o-bonus/stochastic/post_reward/paired_p:3}$); a plausible reason, which we did not test, is that a transition looks stale only because it is noisy and the bonus drives pointless re-exploration.

\textbf{Bonus weight (Table~\ref{tab:kappa}).} The default $\kappa=10^{-3}$ was fixed in advance and not tuned. At $\kappa\ge10^{-2}$ Dyna-Q+ collapses on every task to a handful of goals: with $n=10$ the planning bonus $\kappa\sqrt{\tau}$ for rarely tried pairs outweighs the unit goal reward, and the agent keeps chasing it. Smaller $\kappa=10^{-4}$ retains the static and stochastic performance but loses most of the post-change advantage (the shortcut gain disappears; blocking falls between Dyna-Q and the default). The benefit therefore lives in a narrow band of $\kappa$ that depends on the horizon and reward scale, and we only sampled four values. The $n$ sweep (Table~\ref{tab:sweepn}, Fig.~\ref{fig:sweepn}) is the second sensitivity analysis: it is where the planning-step trade-off of the title is read off.
